% Canonical LaTeX source; imported and revised from the legacy Markdown.
\chapter{Successioni, serie di funzioni e serie di potenze}
\label{chap:18}

\begin{obiettivocapitolo}{gestire convergenza uniforme, scambi di limite e serie di potenze}
\prioritamedia o \estensione, secondo il corso. La convergenza puntuale si verifica punto per punto; la convergenza uniforme richiede una stima indipendente dal punto.
\end{obiettivocapitolo}
\begin{proceduraesame}{serie di potenze}
\begin{enumerate}
\item Applica rapporto o radice ai coefficienti per ottenere il raggio $R$.
\item Concludi convergenza assoluta per $|x-x_0|<R$ e divergenza per $|x-x_0|>R$.
\item Sostituisci separatamente $x=x_0-R$ e $x=x_0+R$ e studia le due serie numeriche risultanti.
\item Scrivi l'intervallo con estremi inclusi o esclusi secondo i risultati.
\end{enumerate}
\end{proceduraesame}
\begin{sceltametodo}{convergenza uniforme}
Cerca $\sup_{x\in E}|f_n(x)-f(x)|\to0$ oppure usa il criterio di Weierstrass $|f_n(x)|\le M_n$ con $\sum M_n$ convergente. Su ogni compatto interno all'intervallo di convergenza, una serie di potenze converge uniformemente.
\end{sceltametodo}
\begin{errorefrequente}{scambi non giustificati}
Limite, integrale e derivata non si scambiano automaticamente con una successione o una serie. Verifica il teorema specifico e le sue ipotesi.
\end{errorefrequente}

Si assume \(\varnothing\ne E\subseteq\mathbb R\) e funzioni a valori in \(\mathbb K\in\{\mathbb R,\mathbb C\}\).

Teoria ed esercizi: \href{https://ingegnerismo.it/matematica/scambio-limite-integrale-derivata-esercizi/}{scambio fra limite, integrale e derivata}.

\section{Convergenza puntuale e uniforme}\label{convergenza-puntuale-e-uniforme}

La differenza è nell'ordine dei quantificatori: nella convergenza puntuale la soglia \(N\) può dipendere da \(x\); in quella uniforme una sola soglia deve funzionare contemporaneamente per ogni \(x\in E\).

Convergenza puntuale:

\[
\displaystyle
f_n\xrightarrow[]{p}f\ \text{su }E
\quad\Longleftrightarrow\quad
\forall x\in E\ \forall\varepsilon>0\ \exists N=N(x,\varepsilon):
\ n\ge N
\Longrightarrow
\lvert f_n(x)-f(x)\rvert<\varepsilon.
\]

Per una funzione \(g:E\to\mathbb K\) poniamo, con valore eventualmente infinito,

\[
\displaystyle
\lVert g\rVert_{\infty,E}
=
\sup_{x\in E}\lvert g(x)\rvert\in[0,+\infty].
\]

Sullo spazio delle funzioni limitate questo valore è finito ed è la norma uniforme.

Convergenza uniforme:

\[
\displaystyle
f_n\xrightarrow[]{u}f\ \text{su }E
\quad\Longleftrightarrow\quad
\lVert f_n-f\rVert_{\infty,E}\to0,
\]

\[
\displaystyle
\forall\varepsilon>0\ \exists N=N(\varepsilon):
\ n\ge N
\Longrightarrow
\lvert f_n(x)-f(x)\rvert<\varepsilon
\quad\forall x\in E.
\]

\[
\displaystyle
f_n\xrightarrow[]{u}f
\quad\Longrightarrow\quad
f_n\xrightarrow[]{p}f.
\]

Negazione dell'uniformità:

\[
\displaystyle
\neg\bigl(f_n\xrightarrow[]{u}f\bigr)
\quad\Longleftrightarrow\quad
\exists\varepsilon_0>0\ \forall N\ \exists n\ge N\ \exists x_n\in E:
\ \lvert f_n(x_n)-f(x_n)\rvert\ge\varepsilon_0.
\]

Criterio uniforme di Cauchy:

\[
\displaystyle
f_n\ \text{converge uniformemente su }E
\]

\[
\displaystyle
\Longleftrightarrow
\forall\varepsilon>0\ \exists N\ \forall m,n:
\ m,n\ge N
\Longrightarrow
\sup_{x\in E}\lvert f_n(x)-f_m(x)\rvert<\varepsilon.
\]

Se \(K\ne\varnothing\) è compatto, lo spazio \(C(K)\) con la norma uniforme è completo:

\[
(f_n)\text{ di Cauchy in }\lVert\cdot\rVert_{\infty,K}
\Longrightarrow
\exists f\in C(K):\ \lVert f_n-f\rVert_{\infty,K}\to0.
\]

Teoria: \href{https://ingegnerismo.it/matematica/convergenza-uniforme/}{convergenza uniforme}.

\section{Teoremi di scambio}\label{teoremi-di-scambio}

Continuità:

\[
\displaystyle
f_n\in C(E),\quad
f_n\xrightarrow[]{u}f
\quad\Longrightarrow\quad
f\in C(E).
\]

Scambio dei limiti, con \(x_0\) punto di accumulazione di \(E\) e con ogni \(\ell_n\) esistente e finito:

\[
\displaystyle
f_n\xrightarrow[]{u}f,\qquad
\ell_n=\lim_{\substack{x\to x_0\\x\in E}}f_n(x)
\]

\[
\displaystyle
\Longrightarrow
\ell_n\text{ converge},
\qquad
\lim_{\substack{x\to x_0\\x\in E}}f(x)
=
\lim_{n\to\infty}\ell_n.
\]

Teorema di Dini, nel caso reale:

\[
\displaystyle
E\ \text{compatto},\quad
f_n,f\in C(E),\quad
f_n\xrightarrow[]{p}f,
\]

\[
\displaystyle
f_n\le f_{n+1}\ \ \forall n
\quad\text{oppure}\quad
f_n\ge f_{n+1}\ \ \forall n
\]

\[
\displaystyle
\Longrightarrow
f_n\xrightarrow[]{u}f.
\]

Scambio con l'integrale:

\[
\displaystyle
f_n\in\mathcal R([a,b]),\quad
f_n\xrightarrow[]{u}f
\]

\[
\displaystyle
\Longrightarrow
f\in\mathcal R([a,b]),
\qquad
\lim_{n\to\infty}\int_a^b f_n(x)\,dx
=
\int_a^b f(x)\,dx.
\]

Stima:

\[
\displaystyle
\left\lvert\int_a^b(f_n-f)\right\rvert
\le
(b-a)\lVert f_n-f\rVert_{\infty,[a,b]}.
\]

Scambio con la derivata:

\[
\displaystyle
f_n\in C^1([a,b]),\quad
f_n(x_\ast)\ \text{convergente per un }x_\ast\in[a,b],
\quad
f_n'\xrightarrow[]{u}g
\]

\[
\displaystyle
\Longrightarrow
\exists f\in C^1([a,b]):
\quad
f_n\xrightarrow[]{u}f,
\qquad
f'=g.
\]

Teoria generale: \href{https://ingegnerismo.it/matematica/convergenza-uniforme/}{convergenza uniforme}.

\section{Serie di funzioni}\label{serie-di-funzioni}

Somme parziali:

\[
\displaystyle
S_N(x)=\sum_{n=0}^{N}f_n(x).
\]

Convergenza puntuale e uniforme:

\[
\displaystyle
\sum_{n=0}^{\infty}f_n(x)=S(x)
\quad\Longleftrightarrow\quad
S_N(x)\to S(x),
\]

\[
\displaystyle
\sum f_n\ \text{converge uniformemente}
\quad\Longleftrightarrow\quad
\lVert S_N-S\rVert_{\infty,E}\to0.
\]

Criterio uniforme di Cauchy:

\[
\displaystyle
\forall\varepsilon>0\ \exists N\ \forall p,q:
\ q\ge p\ge N
\Longrightarrow
\sup_{x\in E}
\left\lvert\sum_{n=p}^{q}f_n(x)\right\rvert<\varepsilon.
\]

Condizione necessaria:

\[
\displaystyle
\sum f_n\ \text{converge uniformemente}
\quad\Longrightarrow\quad
\lVert f_n\rVert_{\infty,E}\to0.
\]

Convergenza totale:

\[
\displaystyle
\sum_{n=0}^{\infty}\lVert f_n\rVert_{\infty,E}<+\infty.
\]

Criterio di Weierstrass:

\[
\displaystyle
M_n\ge0,\qquad\lvert f_n(x)\rvert\le M_n\quad\forall x\in E,
\qquad
\sum M_n<+\infty
\]

\[
\displaystyle
\Longrightarrow
\sum f_n\ \text{converge totalmente, uniformemente e puntualmente in modo assoluto}.
\]

Stima uniforme del resto:

\[
\displaystyle
\sup_{x\in E}
\left\lvert
\sum_{n=N+1}^{\infty}f_n(x)
\right\rvert
\le
\sum_{n=N+1}^{\infty}M_n.
\]

Teoria: \href{https://ingegnerismo.it/matematica/convergenza-uniforme/}{convergenza uniforme} e \href{https://ingegnerismo.it/matematica/convergenza-totale/}{convergenza totale}.

\section{Continuità, integrazione e derivazione delle serie}\label{continuituxe0-integrazione-e-derivazione-delle-serie}

Continuità:

\[
\displaystyle
f_n\in C(E),\quad
\sum f_n\ \text{converge uniformemente}
\quad\Longrightarrow\quad
\sum f_n\in C(E).
\]

Integrazione:

\[
\displaystyle
f_n\in\mathcal R([a,b]),\quad
\sum f_n\ \text{converge uniformemente}
\]

\[
\displaystyle
\Longrightarrow
\int_a^b\sum_{n=0}^{\infty}f_n(x)\,dx
=
\sum_{n=0}^{\infty}\int_a^b f_n(x)\,dx.
\]

Derivazione:

\[
\displaystyle
f_n\in C^1([a,b]),\quad
\sum f_n(x_\ast)\ \text{converge per un }x_\ast\in[a,b],
\quad
\sum f_n'\ \text{converge uniformemente}
\]

\[
\displaystyle
\Longrightarrow
\sum f_n\ \text{converge uniformemente},
\qquad
\left(\sum_{n=0}^{\infty}f_n\right)'
=
\sum_{n=0}^{\infty}f_n'.
\]

Dirichlet uniforme:

\[
\displaystyle
\sup_{\substack{N\ge0\\x\in E}}
\left\lvert\sum_{n=0}^{N}a_n(x)\right\rvert<+\infty,
\quad
b_n(x)\ge0,
\quad
b_{n+1}(x)\le b_n(x),
\quad
\sup_{x\in E}b_n(x)\to0
\]

\[
\displaystyle
\Longrightarrow
\sum a_n(x)b_n(x)\ \text{converge uniformemente}.
\]

Abel uniforme:

\[
\displaystyle
\sum a_n(x)\ \text{converge uniformemente},
\quad
b_n(x)\ \text{monotona in }n\ \text{con verso comune},
\quad
\sup_{\substack{n\ge0\\x\in E}}\lvert b_n(x)\rvert<+\infty
\]

\[
\displaystyle
\Longrightarrow
\sum a_n(x)b_n(x)\ \text{converge uniformemente}.
\]

Teoria: \href{https://ingegnerismo.it/matematica/convergenza-totale/}{convergenza totale} e \href{https://ingegnerismo.it/matematica/criteri-di-abel-e-dirichlet/}{criteri di Abel e Dirichlet}.

\section{Serie di potenze}\label{serie-di-potenze}

Nella serie

\[
\displaystyle
\sum_{n=0}^{\infty}c_n(x-x_0)^n,
\]

\(x_0\) è il centro, \(c_n\) il coefficiente di ordine \(n\) e \(R\) il raggio di convergenza. Il termine con \(n=0\) è il termine costante \(c_0\), anche nel centro \(x=x_0\), secondo la convenzione contestuale illustrata nella guida iniziale. All'interno di \(|x-x_0|<R\) la serie converge assolutamente; i due estremi, quando finiti, vanno studiati separatamente.

Formula di Cauchy--Hadamard:

\[
\displaystyle
\rho=\limsup_{n\to\infty}\sqrt[n]{\lvert c_n\rvert},
\qquad
R=\frac1\rho.
\]

Qui \(R=+\infty\) quando \(\rho=0\) e \(R=0\) quando \(\rho=+\infty\): sono convenzioni per il raggio di convergenza, non divisioni ordinarie per zero o per infinito.

Se \(c_{n+1}\ne0\) definitivamente ed esiste il limite, finito o esteso:

\[
\displaystyle
R=\lim_{n\to\infty}
\left\lvert\frac{c_n}{c_{n+1}}\right\rvert.
\]

Regione di convergenza:

\begin{longtable}[]{@{}
  >{\raggedright\arraybackslash}p{(\columnwidth - 6\tabcolsep) * \real{0.2500}}
  >{\raggedright\arraybackslash}p{(\columnwidth - 6\tabcolsep) * \real{0.2500}}
  >{\raggedright\arraybackslash}p{(\columnwidth - 6\tabcolsep) * \real{0.2500}}
  >{\raggedright\arraybackslash}p{(\columnwidth - 6\tabcolsep) * \real{0.2500}}@{}}
\toprule\noalign{}
\begin{minipage}[b]{\linewidth}\raggedright
Raggio
\end{minipage} & \begin{minipage}[b]{\linewidth}\raggedright
Convergenza
\end{minipage} & \begin{minipage}[b]{\linewidth}\raggedright
Divergenza
\end{minipage} & \begin{minipage}[b]{\linewidth}\raggedright
Bordi
\end{minipage} \\
\midrule\noalign{}
\endhead
\bottomrule\noalign{}
\endlastfoot
\(R=0\) & \(x=x_0\) & \(x\ne x_0\) & nessun bordo da studiare \\
\(0<R<+\infty\) & assoluta per \(\lvert x-x_0\rvert<R\) & per \(\lvert x-x_0\rvert>R\) & \(x_0-R\) e \(x_0+R\), separatamente \\
\(R=+\infty\) & assoluta per ogni \(x\in\mathbb R\) & --- & nessun bordo finito \\
\end{longtable}

Casi estremi:

\[
\displaystyle
\sum_{n=0}^{\infty}n!(x-x_0)^n:
\qquad R=0,
\]

\[
\displaystyle
\sum_{n=0}^{\infty}\frac{(x-x_0)^n}{n!}:
\qquad R=+\infty.
\]

Convergenza totale sui compatti interni, per \(0\le r<R\):

\[
\displaystyle
\sum_{n=0}^{\infty}
\sup_{\lvert x-x_0\rvert\le r}
\lvert c_n(x-x_0)^n\rvert
=
\sum_{n=0}^{\infty}\lvert c_n\rvert r^n
<+\infty.
\]

Comportamenti al bordo per serie centrate in \(0\):

\begin{longtable}[]{@{}
  >{\raggedright\arraybackslash}p{(\columnwidth - 6\tabcolsep) * \real{0.2308}}
  >{\raggedleft\arraybackslash}p{(\columnwidth - 6\tabcolsep) * \real{0.3077}}
  >{\raggedright\arraybackslash}p{(\columnwidth - 6\tabcolsep) * \real{0.2308}}
  >{\raggedright\arraybackslash}p{(\columnwidth - 6\tabcolsep) * \real{0.2308}}@{}}
\toprule\noalign{}
\begin{minipage}[b]{\linewidth}\raggedright
Serie
\end{minipage} & \begin{minipage}[b]{\linewidth}\raggedleft
\(R\)
\end{minipage} & \begin{minipage}[b]{\linewidth}\raggedright
\(x=-R\)
\end{minipage} & \begin{minipage}[b]{\linewidth}\raggedright
\(x=R\)
\end{minipage} \\
\midrule\noalign{}
\endhead
\bottomrule\noalign{}
\endlastfoot
\(\displaystyle\sum_{n=0}^{\infty}x^n\) & \(1\) & diverge & diverge \\
\(\displaystyle\sum_{n=1}^{\infty}\dfrac{x^n}{n}\) & \(1\) & converge condizionatamente & diverge \\
\(\displaystyle\sum_{n=1}^{\infty}\dfrac{x^n}{n^2}\) & \(1\) & converge assolutamente & converge assolutamente \\
\end{longtable}

\section{Operazioni sulle serie di potenze}\label{operazioni-sulle-serie-di-potenze}

Posto

\[
\displaystyle
F(x)=\sum_{n=0}^{\infty}c_n(x-x_0)^n,
\qquad
\lvert x-x_0\rvert<R,
\]

per \(k\in\mathbb N_0\):

\[
\displaystyle
F^{(k)}(x)
=
\sum_{n=k}^{\infty}
\frac{n!}{(n-k)!}\,
c_n(x-x_0)^{n-k},
\]

\[
\displaystyle
c_k=\frac{F^{(k)}(x_0)}{k!}.
\]

Integrazione:

\[
\displaystyle
\int_{x_0}^{x}F(t)\,dt
=
\sum_{n=0}^{\infty}
\frac{c_n}{n+1}(x-x_0)^{n+1}.
\]

Le serie derivata e integrata hanno raggio \(R\).

Somma e prodotto, per \(\lvert x-x_0\rvert<\min\{R_A,R_B\}\):

\[
\displaystyle
\sum a_n(x-x_0)^n+\sum b_n(x-x_0)^n
=
\sum(a_n+b_n)(x-x_0)^n,
\]

\[
\displaystyle
\left(\sum a_n(x-x_0)^n\right)
\left(\sum b_n(x-x_0)^n\right)
=
\sum_{n=0}^{\infty}
\left(\sum_{k=0}^{n}a_kb_{n-k}\right)
(x-x_0)^n.
\]

Teorema di Abel al bordo destro, per \(0<R<+\infty\):

\[
\displaystyle
\sum_{n=0}^{\infty}c_nR^n\ \text{converge}
\quad\Longrightarrow\quad
\lim_{x\to(x_0+R)^-}F(x)
=
\sum_{n=0}^{\infty}c_nR^n.
\]

Bordo sinistro:

\[
\displaystyle
\sum_{n=0}^{\infty}c_n(-R)^n\ \text{converge}
\quad\Longrightarrow\quad
\lim_{x\to(x_0-R)^+}F(x)
=
\sum_{n=0}^{\infty}c_n(-R)^n.
\]

Teoria, bordi e operazioni: \href{https://ingegnerismo.it/matematica/serie-di-potenze/}{serie di potenze}.

\section{Serie di potenze notevoli}\label{serie-di-potenze-notevoli}

Identità fondamentali:

\[
\sum_{n=0}^{\infty}x^n=\frac1{1-x},
\qquad |x|<1,
\]

\[
e^x=\sum_{n=0}^{\infty}\frac{x^n}{n!},
\qquad x\in\mathbb R,
\]

\[
\sin x=\sum_{n=0}^{\infty}(-1)^n\frac{x^{2n+1}}{(2n+1)!},
\qquad
\cos x=\sum_{n=0}^{\infty}(-1)^n\frac{x^{2n}}{(2n)!},
\qquad x\in\mathbb R,
\]

\[
-\ln(1-x)=\sum_{n=1}^{\infty}\frac{x^n}{n},
\qquad -1\le x<1,
\]

\[
\ln(1+x)=\sum_{n=1}^{\infty}(-1)^{n-1}\frac{x^n}{n},
\qquad -1<x\le1,
\]

\[
\arctan x=\sum_{n=0}^{\infty}(-1)^n\frac{x^{2n+1}}{2n+1},
\qquad |x|\le1.
\]

\[
\arcsin x
=\sum_{n=0}^{\infty}
\frac{\binom{2n}{n}}{4^n(2n+1)}x^{2n+1},
\qquad |x|\le1.
\]

Per \(\alpha\in\mathbb R\):

\[
(1+x)^\alpha
=\sum_{n=0}^{\infty}\binom{\alpha}{n}x^n,
\qquad |x|<1,
\]

con gli estremi da discutere in base ad \(\alpha\).

Agli estremi del raggio di convergenza il comportamento va sempre controllato separatamente: la serie può convergere assolutamente, convergere soltanto condizionatamente oppure divergere. Nei casi elencati sopra:

\begin{itemize}
\tightlist
\item
  \(\sum_{n=0}^{\infty}x^n\) diverge sia per \(x=1\) sia per \(x=-1\);
\item
  la serie di \(-\ln(1-x)\) converge condizionatamente per \(x=-1\) e diverge per \(x=1\);
\item
  la serie di \(\ln(1+x)\) diverge per \(x=-1\) e converge condizionatamente per \(x=1\);
\item
  la serie di \(\arctan x\) converge condizionatamente per \(x=1\) e per \(x=-1\);
\item
  la serie di \(\arcsin x\) converge assolutamente per \(x=1\) e per \(x=-1\);
\item
  per la serie binomiale il comportamento dipende da \(\alpha\) e va studiato separatamente nei due estremi.
\end{itemize}
